\section{Three Connected Questions}
\textbf{Economics:} Under quality uncertainty, does a signaling mechanism narrow adverse selection enough to sustain stranger-to-stranger exchange, or does it merely delay market unraveling? \textbf{Computer Science:} Can a lightweight, client-side AI signal make item condition inspectable without a dedicated compute runtime or opaque large-model judgment? \textbf{Behavioral Science:} Does opponent-specific reciprocity emerge among strangers under minimal observability, and how much does an external signal change that emergence? \textbf{Integrated question:} Under a fixed compute and privacy budget, does an AI-assisted trust signal increase completed stranger-to-stranger exchanges on a campus circular-economy platform, and is its effect closer to Akerlof's ``not enough to prevent collapse'' or Leung et al.'s ``enough to bootstrap cooperation''?

\section{Answer to Q1: Economics}
Akerlof's~\cite{akerlof1970} classical result shows that unresolved information asymmetry about quality can depress or collapse a market. A campus exchange platform faces exactly this problem: buyers cannot verify a listed item's true condition before a transaction. Players are the requester (buyer) and the counterpart (seller); actions are list, verify, and trust/decline; payoffs depend on whether the transaction completes and whether the item matches its stated condition. The proposed mechanism narrows, but does not eliminate, this asymmetry through a lightweight verification signal, tested against a no-signal baseline to measure how much of the adverse-selection gap the signal actually closes.

\section{Answer to Q2: Computation}
The mechanism is implemented as a Hugging Face Static Space (\DemoLink), using client-side HTML/CSS/JavaScript only, with no dedicated CPU/GPU runtime, backend, database, or API key. The prototype runs a ten-round trust game under two conditions -- AI-Verified Badge ($p=0.75$) and No Badge/Manual Review ($p=0.45$) -- with seeded logic (seed = 42) so every reported outcome is reproducible. Running the Space's own seeded logic for one ten-round session per condition returned reciprocation in 8 of 10 rounds (80\%) under the Badge condition and 3 of 10 rounds (30\%) under No Badge; these are the artifact's own encoded, illustrative probabilities, not measurements of real campus users. A companion Colab notebook (\ColabLink) reimplements the same probability model in Python to independently verify these round-by-round outcomes.

\clearpage\balance
\section{Answer to Q3: Behavior}
Morozov and Feigel~\cite{morozovfeigel2026} show that strangers can sustain cooperation without full observability through opponent-specific responses, rather than requiring repeated identity tracking. The competing explanation is that any observed cooperation gap reflects only the stated badge probability, not a genuine behavioral response to trust signaling. Discriminating evidence would be a real-user trial where completion rates diverge from the encoded probabilities in a direction opponent-specific reciprocity predicts. Synthetic results from the seeded prototype cannot themselves establish human behavior; a feasible next test is a small live pilot comparing self-reported trust with actual completion under each condition, and a null result -- no behavioral difference beyond the encoded probability -- would change this conclusion.

\section{Advanced Development and Future Directions}
This proposal builds on Akerlof's~\cite{akerlof1970} 1970 Nobel-recognized foundation on information asymmetry in markets, and on the algorithmic-game-theory lineage connecting cooperation, minimal observability, and learning~\cite{leungetal2026}. The enduring human question is how strangers establish trust absent repeated identity or full observability. Today's distinctive environment is a resource-constrained campus platform where verification cannot rely on costly compute or persistent surveillance; current research leaves open how much a lightweight, inspectable signal actually substitutes for full observability, rather than assuming it fully solves adverse selection.

\begin{table}[h]\caption{From laureate foundations to a new interdisciplinary contribution.}
\begin{tabularx}{\columnwidth}{@{}p{1.6cm}Y@{}}\toprule
\textbf{Horizon}&\textbf{Milestone and evidence}\\\midrule
Foundation&Akerlof (1970): adverse selection under quality uncertainty; Leung et al. (2026, AAAI): cooperative learning under minimal observability.\\
PS1&Ten-round trust game (Badge vs.\ No Badge), seeded logic, Colab-verified outputs (8/10 vs.\ 3/10).\\
Next study&Add a live campus pilot comparing self-reported trust with actual completion under each condition.\\
Longer term&Auditable accountability log (Haztrak-style) piloted with the campus sustainability office.\\
2056&A lightweight, resource-aware trust-signaling protocol adopted by resource-constrained institutions.\\\bottomrule
\end{tabularx}\end{table}

\noindent\textbf{Expected 2056 Nobel Prize and Turing Award contribution statement.}
By 2056, I aim to contribute a lightweight, auditable trust-signaling protocol for stranger-to-stranger circular-economy platforms by integrating opponent-specific reciprocity under minimal observability~\cite{morozovfeigel2026} (behavioral advance), signaling mechanisms that counteract adverse selection under quality uncertainty~\cite{akerlof1970} (economic advance), and resource-aware, inspectable AI verification components rather than opaque large-model judgments~\cite{leungetal2026} (computational advance), so that resource-constrained institutions can adopt trust infrastructure without unaccountable AI or unsustainable compute cost. This is an aspiration supported by course-scale milestones, not a prediction of an award.

\subsection*{Open Science Statement}
\textbf{GitHub:} \GitHubLink\\
\textbf{Google Colab:} \ColabLink\par
The repository contains the deployed Space's HTML/CSS/JS source, a Python Colab notebook reimplementing the seeded trust-game logic (seed = 42), and this proposal's LaTeX source. Dependencies are limited to standard Python scientific packages installed in the notebook's setup cell. Shortest reproducible run: open the Colab link and run all cells to regenerate the 8/10 and 3/10 outcomes reported above. Released materials include the full game logic and seeded outputs; validation with real campus users remains a benefit awaiting evaluation, not a released result.

\subsection*{Statement of Contribution to the UN's SDGs}
\textbf{Hugging Face:} \DemoLink\\
Through this open-access interactive trust game, this project aims to support \textbf{SDG Target 12.5}~\cite{sdg12} (substantially reduce waste generation through prevention, reduction, recycling, and reuse) by letting students and staff directly experience how a trust signal changes cooperation incentives in a stranger-to-stranger exchange. Mechanism: an AI-assisted trust signal lowers information asymmetry between strangers, raising the completed exchange/donation rate and diverting items from the waste stream. Observable indicator: the share of initiated exchange/donation requests that reach completion on the campus platform. Distinguishing released from evaluated benefits: the demo itself is released and playable now; its effect on real campus exchange behavior remains to be evaluated.

\subsection*{Submission milestones}
First draft: September 13, 2026, 11:00 P.M. Beijing time (UTC+8). Class review: September 14. Proposed: rebuttal September 16 and revised submission September 20, both at 11:00 P.M. Update Appendix~\ref{app:abstract} with each reflection and revision.
